\chapter{Native Python, Zero-Change Checking, and Mixed Mode}

One of the repository's recurring practical commitments is that users should not be forced
into a theorem-prover syntax cliff before they receive useful results. This commitment has
two implementation faces: zero-change checking for existing code and mixed-mode syntax for
new code.

\section*{Zero-change verification}
The package-level quick APIs and zero-change tools are designed so ordinary Python can be
analyzed without being rewritten first. The repository's exports around
\py{hybrid_check()}, \py{hybrid_analyze()}, and implicit presheaf extraction form a gentle
entry point into the larger architecture. This is important because it shows that the
hybrid system is not reserved for carefully annotated greenfield projects.

The right way to read zero-change verification is as inference of local sections from
existing program structure. The analyzer extracts sites from source, docstrings, control
flow, and library use, then runs structural and semantic passes over those inferred sites.
Trust labels remain visible, so the user can distinguish inferred, solver-backed, and
residual facts.

\section*{Mixed-mode syntax}
The mixed-mode syntax in \repo{src/deppy/hybrid/mixed_mode/syntax.py} is the converse
adoption path. Instead of starting from bare existing code, one starts from ordinary Python
augmented with controlled surface forms such as \py{hole()}, \py{spec()},
\py{@guarantee}, \py{assume()}, and \py{check()}. The crucial invariant named in the file
is graceful degradation: before compilation the forms either pass through, return safe
defaults, or raise a controlled placeholder exception. This allows editors, tests, and
ordinary Python execution to remain useful.

\begin{lstlisting}
from deppy.hybrid import guarantee, assume, check, hole

@guarantee("returns a sorted list with no duplicates")
def unique_sorted(xs: list[int]) -> list[int]:
    assume("xs is finite")
    ys = hole("remove duplicates from xs")
    check("ys preserves the elements of xs")
    return sorted(ys)
\end{lstlisting}

The compiler path for these forms reveals the deep relation between the plain and hybrid
systems. Surface forms compile into plain obligations such as refinements and Sigma-like
witness objects, then re-enter the hybrid pipeline as trust-bearing contracts and proof
obligations. This makes mixed mode the public handshake between the two type systems.

\section*{Why this matters for the paper's title}
If cohomological typing is really ``for absolutely everything,'' it must include both
legacy programs and new proof-oriented code. Zero-change checking handles the first case;
mixed mode handles the second. Together they show that the architecture is not only a
research kernel but also a migration strategy.
